
{\actuality}
Искусственный интеллект (ИИ) всё шире применяется в робототехнических задачах, где роботу необходимо принимать решения в процессе взаимодействия со средой. Обычно обучение таких систем проводится на заранее собранных данных или в симуляторах. Однако при переносе обученной стратегии (policy) в реальные условия робот может столкнуться с ситуациями, которые не были учтены заранее, либо с проблемой переноса из симуляции в реальность (sim-to-real). Решением обеих проблем является дообучение робота непосредственно в реальных условиях~\citeinsynopsis{jiang2025irl}.

Методы обучения с подкреплением (Reinforcement Learning, RL) являются одним из перспективных подходов к такому дообучению~\citeinsynopsis{polydoros2017survey,moerland2023model}. В этих методах агент (робот) обучается методом проб и ошибок: он получает вознаграждение за действия, которые приближают его к желаемому поведению. Важной частью обучения является исследование среды, в ходе которого агенту необходимо обнаружить оптимальную последовательность действий, а затем использовать её для решения поставленной задачи. Однако в реальных условиях исследование не может быть произвольным: действия робота могут нанести вред людям, окружающей среде или самой робототехнической системе. Поэтому агент должен не только исследовать среду в поиске полезных действий, но и делать это безопасно. Эту задачу рассматривает область безопасного обучения с подкреплением (Safe Reinforcement Learning, Safe RL)~\citeinsynopsis{altman1999constrained,gu2022review}.

В безопасном обучении с подкреплением, помимо вознаграждения, вводится штраф, который численно описывает нарушения ограничений безопасности~\citeinsynopsis{altman1999constrained}. В результате агент должен решать двойственную задачу: достигать цели и одновременно соблюдать ограничения. Например, робот должен не только прийти в заданную точку, но и не сталкиваться с препятствиями, не превышать скорость или снижать её рядом с людьми.

Существующие подходы решают эту проблему по-разному. Методы на основе функции лагранжиана применимы к широкому классу сред, но обычно не обеспечивают строгих гарантий безопасности~\citeinsynopsis{stooke2020pidlag,ray2019benchmarking,yang2021wcsac}. Методы на основе барьерных функций и других механизмов фильтрации действий могут давать более строгий контроль поведения, однако часто требуют знания модели объекта управления, явного вида функции штрафов и её дифференцируемости~\citeinsynopsis{dalal2018safe,sac-rcbf,koller2018learning,hunt2021verifiably,yangsafe}. Эти требования ограничивают их применение в сложных робототехнических задачах, где полная информация о среде и точная модель динамики часто недоступны.

Одним из способов уменьшить число небезопасных взаимодействий является использование моделей мира~\citeinsynopsis{polydoros2017survey,moerland2023model}. В таких подходах, помимо стратегии агента, обучается модель, предсказывающая будущие состояния среды при выполнении заданных действий, что позволяет генерировать дополнительный синтетический опыт, а также осуществлять планирование в рамках модели мира. Однако безопасное исследование в таких методах зависит от того, насколько корректно модель мира оценивает безопасность состояний~\citeinsynopsis{wabersich2021predictive,jayant2022model,as2022lambda,hogewind2022safe,safedreamer}.

Дополнительные трудности возникают в задачах с большим горизонтом, которые характеризуются разреженным вознаграждением, высокой размерностью пространства состояний и действий, а также большим количеством шагов. Соответственно, агенту необходимо выполнять длинную последовательность действий, чтобы получить полезный сигнал для обучения. Одним из способов справиться с этой проблемой является иерархическое обучение с подкреплением~\citeinsynopsis{pateria2021hrl,sutton1999options,bacon2017optioncritic}. В таких методах верхнеуровневая стратегия выбирает промежуточные подцели или навыки, а низкоуровневая стратегия отвечает за их достижение~\citeinsynopsis{levy2017hierarchical,nachum2018data,hrac}. Это позволяет разбить исходную задачу на более короткие подзадачи и направить исследование агента к достижению конечной цели~\citeinsynopsis{safetylayerhiro,lyapunovrrt,iahrl}.

Объединение моделей мира и иерархического обучения с подкреплением позволяет использовать оба преимущества одновременно: декомпозицию поведения на промежуточные подцели и предварительную оценку риска их достижения~\citeinsynopsis{gorbov2026ites}. Такой подход особенно важен в навигационных и робототехнических задачах, где одноуровневые безопасные стратегии часто становятся чрезмерно консервативными или не достигают цели.

При переходе от симуляционных задач к взаимодействию с человеком возникает ещё одна проблема: цели и правила поведения часто задаются не в виде явной функции штрафа, а на естественном языке~\citeinsynopsis{andreas2015alignment,anderson2018vision,misra2018mapping,hermann2020directions,hao2020pretraining,chen2019touchdown,de_vries2018talk}. В этом случае агенту необходимо связать текстовое ограничение с объектами, состояниями среды и допустимыми последовательностями действий. Поэтому методы, ориентированные только на заранее заданные численные штрафы, не покрывают весь класс задач, возникающих при использовании автономных агентов в человеко-ориентированных средах. Для проверки таких методов требуются среды, где языковая инструкция влияет не только на цель, но и на допустимость поведения агента~\citeinsynopsis{hazardworld,prakash2020guiding}.

Таким образом, актуальность данной работы определяется тремя связанными проблемами. Во-первых, в задачах безопасного обучения важно заранее оценивать долгосрочные последствия действий: короткий горизонт планирования может не выявлять будущие нарушения безопасности, а ошибки в оценке нарушений приводят либо к небезопасным траекториям, либо к чрезмерно осторожному поведению. Во-вторых, безопасная стратегия должна сохранять способность к исследованию среды в постановках с большим горизонтом, а не сводить поведение к избеганию любых потенциально рискованных действий. В-третьих, для более реалистичных постановок необходимо учитывать ограничения, заданные не только численно, но и на естественном языке.

{\aim} данной работы является исследование и разработка методов безопасного обучения с подкреплением, которые используют модели мира и иерархическую структуру поведения для повышения эффективности исследования среды при соблюдении ограничений безопасности в робототехнических задачах.

{\tasksintro}
\begin{enumerate}[beginpenalty=10000]
  \item Провести анализ существующих методов безопасного обучения с подкреплением и сред, используемых для их экспериментальной оценки.
  \item Разработать подход для учёта долгосрочного риска в методах безопасного обучения с подкреплением на основе модели мира.
  \item Разработать метод иерархического обучения с подкреплением, направленный на уменьшение консервативности безопасной стратегии и повышение её способности к исследованию среды.
  \item Разработать метод генерации безопасных подцелей с использованием модели мира для обеспечения безопасности как у верхнеуровневой стратегии, так и у низкоуровневой для решения робототехнических задач с большим горизонтом.
  \item Разработать бенчмарк для обучения методов безопасного обучения с подкреплением и для оценки их способностей к пониманию задачи и ограничений, заданных при помощи естественного языка.
\end{enumerate}

{\novelty}
\begin{enumerate}[beginpenalty=10000]
  \item В рамках метода SPOWL (Safe Planning and Policy Optimization via World Model Learning, безопасное планирование и оптимизация стратегии с использованием модели мира) предложена функция оценки будущих накопленных штрафов в рамках модели мира для улучшения скрытых представлений, а также для более точной оценки безопасности траекторий при планировании. Эксперименты на средах Safety Gym показали лучшую производительность и безопасность по сравнению с предыдущими методами безопасного обучения с подкреплением.
  \item Разработан метод SG-Safe (Enhancing Safe Exploration Through Subgoal Guidance) для улучшения исследовательских способностей безопасных стратегий за счёт направления исследования через промежуточные подцели. На этапе исполнения используется только безопасная стратегия. Метод показывает лучшее качество в среде POLAMP в сравнении с методами безопасного обучения с подкреплением и методами на основе классического планирования.
  \item Разработан метод ITES (Imagine to Ensure Safety in Hierarchical Reinforcement Learning) для генерации безопасных подцелей в задачах с большим горизонтом, использующий модель мира для обеспечения безопасного достижения этих подцелей. Метод достигает лучшего качества в задачах навигации и манипуляции с большим горизонтом в сравнении с методами на основе моделей мира в безопасном обучении с подкреплением.
  \item Разработан бенчмарк Caged CrafText (Caged CrafText: Language-Grounded Safety Constraints for Multimodal Reinforcement Learning in CrafText) для оценки способностей к пониманию языка, связывания языка с объектами из среды, памяти, безопасного исследования. Предложенный бенчмарк позволяет как обучать, так и оценивать методы безопасного обучения с подкреплением в сложной динамической среде.
\end{enumerate}

{\theory}
В работе показано, что в задачах с большим горизонтом у нынешних методов безопасного обучения с подкреплением ухудшается способность к исследованию, и что такая проблема решается с помощью иерархического обучения с подкреплением в методах SG-Safe и ITES.

Показано, что безопасное исследование в методах обучения с подкреплением на основе модели мира зависит от корректной оценки моделью мира нарушений на сгенерированных траекториях. Недооценка будущих штрафов приводит к небезопасному исследованию, а переоценка~--- к чрезмерно консервативному поведению агента.

Caged CrafText позволяет обучать и оценивать методы безопасного обучения с подкреплением в задачах, где цель и ограничения заданы на естественном языке. Разделение задач на бюджетные, последовательные, математические и отношенческие ограничения позволяет оценивать разные аспекты поведения агента: понимание инструкции, связывание текста с объектами и состояниями среды, память, долгосрочное планирование и безопасное исследование.

{\influence} заключается в возможности применения разработанных методов для робототехнических платформ, где требуется заданный уровень безопасности и производительности, особенно для задач, требующих большого количества действий от робота, в том числе для адаптивного обучения робота в реальных условиях.

Компоненты метода SPOWL, отвечающие за оценку будущих накопленных штрафов действий, могут быть применены в робототехнических задачах безопасного исследования, где агент должен заранее оценивать риск кандидатных действий до их выполнения в реальной среде. Это позволяет уменьшить число небезопасных действий при обучении и одновременно избежать чрезмерно консервативного поведения агента.

Разработанные методы SG-Safe и ITES, направленные на улучшение исследовательских способностей безопасной стратегии при помощи иерархического обучения, могут быть применены для адаптации мобильных роботов к новым условиям эксплуатации, например в задачах автономной инспекции промышленных объектов, патрулирования охраняемых территорий, обследования аварийных зон, навигации в шахтах и других труднодоступных пространствах, где требуется исследовать среду при ограничении числа небезопасных действий.

{\methods}
Для оптимизации стратегий использовались методы актор-критик, включая метод оптимизации ближайшей стратегии (Proximal Policy Optimization, PPO) и метод мягкого актор-критика (Soft Actor-Critic, SAC), а также методы на основе функции лагранжиана для учёта ограничений безопасности. Для формализации задач применялись марковские процессы принятия решений с ограничениями (CMDP). В качестве одного из базовых методов сравнения рассматривался SafeDreamer, основанный на архитектуре DreamerV3 и использующий модель мира для планирования с учётом ограничений безопасности. Для оценки предложенных методов проводились численные эксперименты в симуляционных средах безопасного обучения с подкреплением. Основными метриками выступали эпизодическое вознаграждение (Reward), эпизодические штрафы (Cost), успешность выполнения задачи (Success Rate). Программная реализация выполнялась на языке Python~3 с использованием JAX и PyTorch. В качестве основного программного интерфейса для симуляционных сред использовались Gymnasium и Safety Gymnasium, а также собственные реализации сред для отдельных экспериментов. Эксперименты проводились на графических ускорителях NVIDIA RTX-серии и NVIDIA H100. Логирование результатов экспериментов выполнялось с использованием Comet ML, а визуализация и анализ результатов~--- с использованием Matplotlib и стандартных инструментов обработки численных данных. Репозитории с открытым кодом: SPOWL\footnote{\url{https://github.com/AIRI-Institute/SPOWL}}, SG-Safe\footnote{\url{https://github.com/Gricha1/SPEIS}}, ITES\footnote{\url{https://github.com/Gricha1/ITES_alg}}, CrafText\footnote{\url{https://github.com/AIRI-Institute/CrafText}}, Caged CrafText\footnote{\url{https://github.com/Gricha1/CAGED-CrafText}}.

{\defpositions}
\begin{enumerate}[beginpenalty=10000]
  \item Способ построения и агрегации функций оценки будущих накопленных штрафов в модели мира для безопасного исследования в рамках модели SPOWL. Показано, что устойчивое безопасное исследование в обучении с подкреплением на основе модели мира требует оценки не только мгновенных штрафов, но и оценки будущих накопленных штрафов.
  \item Метод SG-Safe для безопасного иерархического обучения с подкреплением, использующий генерацию промежуточных подцелей для улучшения исследовательских способностей безопасной стратегии. Метод позволяет повысить успешность решения навигационных задач с большим горизонтом при сохранении контролируемого уровня нарушений безопасности.
  \item Метод ITES для безопасного исследования с использованием иерархической структуры и модели мира, в котором верхнеуровневая стратегия генерирует безопасные подцели, а низкоуровневая стратегия использует модель мира и модель штрафов для безопасного достижения этих подцелей. Метод позволяет контролировать безопасность как на уровне выбора подцелей, так и на уровне исполнения действий.
  \item Бенчмарк Caged CrafText для оценки и обучения методов безопасного обучения с подкреплением с естественно-языковыми целями и ограничениями, позволяющий анализировать способность агентов понимать задачи и ограничения, заданные на естественном языке, и соблюдать их в сложной динамической среде.
\end{enumerate}

{\reliability} результатов обеспечивается многократным проведением экспериментов, использованием общепринятых симуляционных сред и сравнением с современными базовыми методами. Достоверность подтверждается также абляционными исследованиями, апробацией результатов на научных семинарах и конференциях, а также публикациями в рецензируемых научных изданиях.

{\probation}
Основные результаты диссертационной работы докладывались и обсуждались на:
\begin{itemize}
  \item семинарах Центра когнитивного моделирования МФТИ (2023, 2024, 2025);
  \item международной конференции European Conference on Artificial Intelligence, ECAI 2025;
  \item конференции Fall into ML 2025;
  \item междисциплинарной конференции по искусственному интеллекту «ИИ-АЗАМАН 2025».
\end{itemize}
Полученные результаты и знания использовались в преподавании курсов по искусственному интеллекту:
\begin{itemize}
  \item курс «Продвинутые методы машинного обучения с подкреплением» (МФТИ, 2026): разработка и проведение практических занятий;
  \item курс «Обучение с подкреплением» (Центр «ПУСК», 2025): разработка и проведение практических и лекционных занятий;
  \item курс «Обучение с подкреплением (углубленный курс)» (ФКН ВШЭ, 2026): проведение лекции и семинара по безопасному обучению с подкреплением.
\end{itemize}

{\publications}
Основные результаты по теме диссертации изложены в 5~публикациях: 2~--- в журналах категории~К1 собственного перечня журналов МФТИ\ifsynopsis\else~\cite{gorbov2026sgsafe,gorbov2026caged}\fi, 2~--- в сборниках трудов конференций уровня~A и~A*\ifsynopsis\else~\cite{latyshev2025spowl,volovikova2025craftext}\fi, 1~--- в виде препринта на arXiv\ifsynopsis\else~\cite{gorbov2026ites}\fi.

{\contribution}
В работе по SPOWL~\citeinsynopsis{latyshev2025spowl} автор предложил функцию оценки будущих накопленных штрафов в рамках модели мира для улучшения скрытых представлений и более точной оценки безопасности траекторий при планировании, а также способ агрегации её в ансамбле. В частности, автор провёл экспериментальный анализ способов агрегации ансамбля функций оценки будущих накопленных штрафов и влияния размера ансамбля на баланс между эффективностью и безопасностью агента. Показано, что корректная ансамблевая оценка будущих накопленных штрафов является необходимым условием безопасного исследования: недооценка риска приводит к нарушениям ограничений, а переоценка риска~--- к консервативному поведению и снижению целевой эффективности.

В работе по SG-Safe~\citeinsynopsis{gorbov2026sgsafe} автор разработал метод безопасного иерархического обучения с подкреплением, направленный на улучшение исследовательских способностей безопасной стратегии за счёт генерации промежуточных подцелей, а также провёл экспериментальную проверку метода на навигационных задачах с большим горизонтом.

В работе по ITES~\citeinsynopsis{gorbov2026ites} автор разработал метод, использующий модель мира и модель штрафов для генерации безопасных подцелей и безопасного достижения этих подцелей низкоуровневой стратегией.

В работе по CrafText Benchmark~\citeinsynopsis{volovikova2025craftext} автор участвовал в разработке задач для оценки способности агента следовать естественно-языковым инструкциям в сложной динамической среде. В частности, были сформированы навигационные и пространственные задачи, связанные с локализацией объектов, выбором направления движения и размещением объектов относительно элементов карты.

В работе~\citeinsynopsis{gorbov2026caged} разработан и протестирован бенчмарк Caged CrafText для оценки агентов безопасного обучения с подкреплением в задачах с естественно-языковыми целями и ограничениями. Сформированы типы ограничений, наборы задач для обучения и тестирования, а также протокол оценки агентов. Проведено экспериментальное сравнение базовых алгоритмов и проанализированы их ограничения при одновременном выполнении языковой инструкции и соблюдении текстового ограничения безопасности в динамической среде.

Во всех указанных работах автор участвовал в обсуждении экспериментальных результатов, их интерпретации и формулировке выводов.

\textbf{Содержание диссертации соответствует паспорту научной специальности 1.2.1.~Искусственный интеллект и машинное обучение, в частности, пунктам:}
\begin{itemize}
  \item \textbf{п.~5}~--- Методы и технологии поиска, приобретения и использования знаний и закономерностей, в том числе~--- эмпирических, в системах искусственного интеллекта. Исследования в области совместного применения методов машинного обучения и классического математического моделирования. Методы и средства использования экспертных знаний.
  \item \textbf{п.~6}~--- Формализация и постановка задач управления и (поддержки) принятия решений на основе систем искусственного интеллекта и машинного обучения. Разработка систем управления с использованием систем искусственного интеллекта и методов машинного обучения в том числе~--- управления роботами, автомобилями, БПЛА и т.п.
  \item \textbf{п.~7}~--- Разработка специализированного математического, алгоритмического и программного обеспечения систем искусственного интеллекта и машинного обучения. Методы и средства взаимодействия систем искусственного интеллекта с другими системами и человеком-оператором.
  \item \textbf{п.~12}~--- Исследования в области «доверенных» систем класса ИИ, включая проблемы формирования тестовых выборок прецедентов, надёжности, устойчивости, переобучения и т.д.
\end{itemize}
